\paragraph{Theorem F.2.13 (Proved: two further rational families excluded integrally).}
Let $C_f,C_g$ denote the degree-four companion matrices.
Take $f=\Phi_5$ and either
\[
 \text{case A: }g=\Phi_{12},\qquad
 \text{case B: }g=\Phi_3\Phi_6.
 \tag{F.36}
\]
Let $W$ be the quotient of
$\bigwedge^2(C_f^{-1},C_g)$ by its invariant line, and
let $V$ be its unique rational extension of F.2.11.
For every invariant integral lattice in $V$ with primitive
maximal square-zero cusp, every simply connected free affine
filling, with multiplicities $m_1,m_2$ divisible by five and
six respectively, fails three-connectivity.
More precisely, put $r=v_3(m_2)\geq1$.
In case A, $\operatorname{Tor}_3 H_3$ has order
$3^{2r-1}$.  In case B,
$\operatorname{Tor}_3 H_2\simeq\mathbb Z/3^r$.
For $m_2=6$ the respective obstructing group is $\mathbb Z/3$.

\begin{proof}
The complete integer representatives are supplied by the
companion and Smith-basis construction in
\texttt{fourfolds\_56\_probe.py} and
\texttt{fourfolds\_56b\_probe.py}.
In each case the source cusp is a rank-one transvection;
the primitive exterior quotient has cusp type $(2,2,1)$,
finite fixed ranks one, and no invariants or coinvariants.
The second finite characteristic polynomial on $V$ is
$(T-1)^2(T+1)^2\Phi_3$ in case A and
$(T-1)^2(T+1)^2\Phi_6$ in case B.
The exact extension equations give an integral primitive
rank-three cusp in both cases.

For completeness these are not merely rational candidates
with impossible twist lattices.  In the displayed integral
bases of the probe, change lattice by the rational basis
matrix $Q$ whose diagonal entries are two and whose only
off-diagonal entries are $Q_{1,4}=Q_{1,6}=1$.
The resulting generators $Q^{-1}AQ,Q^{-1}BQ$ remain
integral and have primitive cusps.
The invariant homology vectors of $\ell$-value one can be
chosen as
\[
\begin{array}{c|cc}
 &u_1&u_2\\\hline
 \mathrm A&(1,3,-3,0,1,0)&(1,-1,1,2,-1,1)\\
 \mathrm B&(1,0,0,1,0,1)&(1,0,0,1,0,0).
\end{array}
\]
Use twists $u_1,-u_2$ and multiplicities five and six.
All nonidentity affine powers are free, and the cusp normal
closure is the saturated $\ker\ell$.
The residual determinant $6-5=1$ proves simple connectivity.
The rational Leray ranks are the same as in F.2.5.

The complete three-adic lattice argument in F.2.12 applies
to each case, not only to these representatives.
The reason is checked by the same exact integer matrices:
the rank-five quotient has no invariants or coinvariants
modulo three, its order-five polynomial is
$(T-1)\Phi_5$ with irreducible $\Phi_5$ modulo three, the
global extension class is primitive, and the rank-six
cusp has unit Smith factors.
Thus each representation has exactly one primitive-cusp
$\mathbb Z_3$-lattice up to homothety.
Both the quotient argument forcing $r_i$ prime to $m_i$
and the finite mapping-torus comparison of F.2.12 apply.

Set $R=\mathbb Z_3$.  All finite cap cohomology groups
through degree four are free over $R$: the nonunit Smith
factors of their vertical exterior differences, when
present, are two.  The global coboundary and coinvariant
Smith factors in degrees one through four likewise have
no factor three.  The rational invariant and parabolic
ranks are those printed in F.2.12.  In particular,
the global invariant groups in degrees two, three and
four vanish, their relevant coinvariants vanish over $R$,
and the degree-two parabolic first cohomology is zero.

For the transformed integral lattices just specified, the
finite degree-two suspension periods at order six are
\[
 \mathrm A:\ (-2,-1,-1),\qquad
 \mathrm B:\ (1/3,0,1/3).
 \tag{F.37}
\]
The order-five periods are integral and its covering degree
is a unit in $R$ for any simply connected multiplicities.
In case B the numerator row at six is $(2,0,2)$; replacing
six by $m_2$ gives the cyclic cokernel $R/3^rR$.
The total-degree-three Leray argument of F.2.12 gives
$H^3(P;R)=R/3^rR$, proving the asserted $H_2$ obstruction.

In case A the obstruction instead occurs in degree four.
The invariant rank-four middle cohomology lattice of the
sixfold torus cover has skew-pairing determinant $6^2$.
This determinant is computed exactly, with the integral
invariant basis, in \texttt{fourfolds\_56a\_newlattice.py}.
The free middle pairing downstairs is unimodular over $R$.
If $I$ is the index of the pullback lattice in the invariant
lattice, pullback degree $m_2$ gives
\[
 I^2\,6^2=m_2^4
 \quad\text{up to a unit of }R.
 \tag{F.38}
\]
Thus the degree-three finite covering cokernel has
three-primary order $3^{2r-1}$.
There are no global invariant degree-three forms to
quotient this cokernel out.  The integral parabolic first
cohomology in degree three is free over $R$, so the torsion
of $H^1(B,R^3g_*R)$ is precisely that finite cokernel.
There are no differentials to or from this middle base
column.  In total degree four the other surviving term
is the free rank-one point-supported cusp class;
the degree-two base term is zero over $R$.
The resulting extension by a free module splits.
Therefore $\operatorname{Tor}_3H^4(P;R)$ has order
$3^{2r-1}$, and the universal coefficient theorem identifies
it with $\operatorname{Tor}_3H_3(P;\mathbb Z)$.
This proves the claimed obstruction in case A.
\end{proof}
